\section*{Capítulo \ref{ch:g11:dissolution} --- A dissolução dos sólidos iônicos e moleculares}

\begin{solution}{exo:g11:dissolution:1}
\ce{NaCl(s) -> Na+(aq) + Cl-(aq)};
\ce{CaCl2(s) -> Ca^{2+}(aq) + 2Cl-(aq)};
\ce{Na2SO4(s) -> 2Na+(aq) + SO4^{2-}(aq)};
\ce{AlCl3(s) -> Al^{3+}(aq) + 3Cl-(aq)}.
\end{solution}

\begin{solution}{exo:g11:dissolution:2}
$[\ce{Na+}] = 2 \times 0.050 = \qty{0.10}{mol/L}$;
$[\ce{SO4^{2-}}] = \qty{0.050}{mol/L}$.
\end{solution}

\begin{solution}{exo:g11:dissolution:3}
A \omterm{def:g11:dissolution:solvation}{hidratação} é o envolvimento do \omterm{def:g9:ions:ion}{íon} por \omterm{def:g7:atoms-and-molecules:molecule}{moléculas} de água presas a ele.
A extremidade do oxigênio ($\delta^-$) aponta para um \omterm{def:g9:ions:ion}{cátion}, a
extremidade dos hidrogênios ($\delta^+$) para um \omterm{def:g9:ions:ion}{ânion}: cargas opostas
se atraem.
\end{solution}

\begin{solution}{exo:g11:dissolution:4}
$M(\ce{AlCl3}) = 27.0 + 3 \times 35.5 = \qty{133.5}{g/mol}$;
$n = 2.67 / 133.5 = \qty{0.0200}{mol}$; $c = 0.0200 / 0.2000 =
\qty{0.100}{mol/L}$; $[\ce{Al^{3+}}] = \qty{0.100}{mol/L}$,
$[\ce{Cl-}] = \qty{0.300}{mol/L}$.
\end{solution}

\begin{solution}{exo:g11:dissolution:5}
A cabeça \ce{-COO-} é iônica, atraída pela água: \omterm{def:g11:dissolution:hydrophilic}{hidrofílica}. A cauda de
17 \omterm{def:g7:atoms-and-molecules:atom}{átomos} de carbono, só com ligações \ce{C-C} e ligações \ce{C-H} quase
apolares, é apolar: \omterm{def:g11:dissolution:hydrophilic}{hidrofóbica}.
\end{solution}

\begin{solution}{exo:g11:dissolution:6}
Sal: água (sólido iônico, \omterm{def:g6:solutions-and-solubility:solute}{solvente} polar). Cera: ciclo-hexano (cadeias
apolares). Açúcar: água (os seus grupos \ce{O-H} formam \omterm{def:g11:polarity-and-cohesion:hydrogen-bond}{ligações de
hidrogênio} com a água). Iodo: ciclo-hexano (\omterm{def:g11:polarity-and-cohesion:polar-molecule}{molécula apolar}).
\end{solution}

\begin{solution}{exo:g11:dissolution:7}
O grupo \ce{O-H} do etanol forma \omterm{def:g11:polarity-and-cohesion:hydrogen-bond}{ligações de hidrogênio} com as \omterm{def:g7:atoms-and-molecules:molecule}{moléculas}
de água, que substituem as que ele rompe: o etanol se encaixa na água. O
hexano, apolar, não consegue formar essas ligações; as \omterm{def:g7:atoms-and-molecules:molecule}{moléculas} de
água, presas umas às outras por \omterm{def:g11:polarity-and-cohesion:hydrogen-bond}{ligações de hidrogênio}, não o deixam
entrar.
\end{solution}

\begin{solution}{exo:g11:dissolution:8}
As caudas \omterm{def:g11:dissolution:hydrophilic}{hidrofóbicas} evitam a água e se dissolvem na gordura, no
centro; as cabeças carregadas são atraídas pela água e ficam do lado de
fora. Todas as \omterm{def:g11:dissolution:amphiphilic}{micelas} têm cargas negativas na superfície, e cargas de
mesmo sinal se repelem: elas ficam separadas.
\end{solution}

\begin{solution}{exo:g11:dissolution:9}
O ciclo-hexano: ele é \omterm{def:g6:solutions-and-solubility:miscible}{imiscível} com a água, por isso forma uma camada
separada que pode ser escoada, e o aroma (apolar) \omterm{def:g2:mixing-and-dissolving:dissolve}{se dissolve} bem nele.
O etanol é \omterm{def:g6:solutions-and-solubility:miscible}{miscível} com a água: não se forma uma segunda camada, e nada
pode ser separado.
\end{solution}

\begin{solution}{exo:g11:dissolution:10}
O ciclo-hexano é menos denso que a água. A camada de baixo, aquosa, é
escoada pela torneira. O iodo está na camada de ciclo-hexano de cima,
que fica no funil.
\end{solution}

\begin{solution}{exo:g11:dissolution:11}
Nada: o sulfato de cobre é iônico e não \omterm{def:g2:mixing-and-dissolving:dissolve}{se dissolve} no ciclo-hexano
apolar. A água continua azul e o ciclo-hexano incolor.
\end{solution}

\begin{solution}{exo:g11:dissolution:12}
$[\ce{Cl-}] = 2c$, logo $c = \qty{0.150}{mol/L}$;
$n = 0.150 \times 0.2500 = \qty{0.0375}{mol}$;
$m = 0.0375 \times 111.1 = \qty{4.17}{g}$.
\end{solution}

\begin{solution}{exo:g11:dissolution:13}
Os \omterm{def:g9:ions:ion}{íons} cálcio e magnésio precipitam os \omterm{def:g9:ions:ion}{íons} estearato como borra: o
sabão é consumido antes de poder formar \omterm{def:g11:dissolution:amphiphilic}{micelas}, e faz pouca espuma.
Só o cálcio: \qty{0.010}{mol} de \ce{Ca^{2+}} toma
\qty{0.020}{mol} de estearato, isto é, $0.020 \times 306.0 =
\qty{6.1}{g}$ de estearato de sódio por litro.
\end{solution}

\begin{solution}{exo:g11:dissolution:14}
Volume \qty{0.2000}{L}. \ce{Na+}: $0.020 / 0.2000 = \qty{0.10}{mol/L}$;
\ce{Ca^{2+}}: $0.010 / 0.2000 = \qty{0.050}{mol/L}$; \ce{Cl-}:
$(0.020 + 0.020) / 0.2000 = \qty{0.20}{mol/L}$. Carga positiva:
$0.10 + 2 \times 0.050 = 0.20$; negativa: $0.20$. Neutra.
\end{solution}

\begin{solution}{exo:g11:dissolution:15}
Um litro de heptano tem massa de \qty{0.68}{kg} e consegue conter
$17 \times 0.68 = \qty{12}{g}$ de iodo, cerca de 40 vezes mais que um
litro de água. Agitados juntos, quase todo o iodo passa para o heptano:
a \omterm{def:g10:chemical-species:extraction}{extração} é eficiente.
\end{solution}

\begin{solution}{pb:g11:dissolution:1}
\textbf{1.} Das rochas por onde a água passou (calcário, dolomita,
gipsita), que se dissolvem um pouco.

\textbf{2.} $[\ce{Ca^{2+}}] = 0.120 / 40.1 = \qty{2.99e-3}{mol/L}$;
$[\ce{Mg^{2+}}] = 0.024 / 24.3 = \qty{9.9e-4}{mol/L}$.

\textbf{3.} \ce{Ca(HCO3)2 -> Ca^{2+}(aq) + 2HCO3-(aq)}; logo
$[\ce{HCO3-}] = 2 \times \num{2.99e-3} = \qty{5.99e-3}{mol/L}$.

\textbf{4.} $\num{2.99e-3} \times 150 = \qty{0.449}{mol}$.

\textbf{5.} $18 \times 12.0 + 35 \times 1.0 + 2 \times 16.0 + 23.0 =
\qty{306.0}{g/mol}$.

\textbf{6.} \ce{C17H35COONa(s) -> C17H35COO-(aq) + Na+(aq)}.

\textbf{7.} Cabeça \ce{-COO-}: \omterm{def:g11:dissolution:hydrophilic}{hidrofílica}; cauda \ce{C17H35-}:
\omterm{def:g11:dissolution:hydrophilic}{hidrofóbica}. Ele tem os dois tipos de parte: é \omterm{def:g11:dissolution:amphiphilic}{anfifílico}.

\textbf{8.} Uma esfera minúscula de \omterm{def:g9:ions:ion}{íons} de sabão, caudas para dentro,
cabeças para fora. As caudas se dissolvem na gordura, que se quebra em
gotículas envolvidas por \omterm{def:g11:dissolution:amphiphilic}{micelas}; as superfícies carregadas mantêm as
gotículas separadas na água, que as leva embora no enxágue.

\textbf{9.} O sabão precisa formar \omterm{def:g9:ions:ion}{íons} e \omterm{def:g11:dissolution:amphiphilic}{micelas}, o que exige água em
volta das cabeças; e a gordura é levada embora pela água do enxágue. No
óleo, a gordura simplesmente continuaria dissolvida na pele.

\textbf{10.} \ce{2C17H35COO- + Ca^{2+} -> Ca(C17H35COO)2}: cargas
$2 \times (-1) + 2 = 0$ à esquerda, $0$ à direita.

\textbf{11.} \ce{C17H35COO-}: excesso $- 2x$; \ce{Ca^{2+}}:
$\num{2.99e-3} - x$; estearato de cálcio: $x$. Os \omterm{def:g9:ions:ion}{íons} cálcio são
limitantes: $x_{\max} = \qty{2.99e-3}{mol}$.

\textbf{12.} $2 \times \num{2.99e-3} = \qty{5.99e-3}{mol}$ por litro.

\textbf{13.} $M = 36 \times 12.0 + 70 \times 1.0 + 4 \times 16.0 + 40.1 =
\qty{606.1}{g/mol}$; $\num{2.99e-3} \times 606.1 = \qty{1.81}{g}$ de
borra por litro.

\textbf{14.} $2 \times \num{9.9e-4} = \qty{1.98e-3}{mol}$ por litro.

\textbf{15.} $2 \times 0.449 = \qty{0.898}{mol}$.

\textbf{16.} $0.898 \times 306.0 = \qty{275}{g}$.

\textbf{17.} Magnésio: $\num{9.9e-4} \times 150 = \qty{0.148}{mol}$,
que toma \qty{0.296}{mol} de estearato, \qty{91}{g} de sabão; ao todo,
cerca de $275 + 91 = \qty{366}{g}$.

\textbf{18.} $275 / 100 \approx 2.7$ barras.

\textbf{19.} O cálcio de um banho desperdiça cerca de \qty{0.27}{kg} de
sabão em borra.
\end{solution}
